% Section 4: Impact
\section{Impact}
\label{sec:impact}

\subsection{Practical value for operations planning}
FHOPS reduces the cost of moving from ad hoc local models to comparable, auditable planning studies. The practical value is not one ``best'' solver run, but a repeatable workflow that keeps inputs, constraints, solver settings, and evaluation outputs aligned across planning cycles.
For planning teams, this supports more defensible schedule revisions because performance differences can be tied to declared assumptions rather than undocumented preprocessing or solver-configuration changes.

\subsection{Research value in the manuscript sequence}
Relative to the review paper \cite{jaffray2025systematic}, this manuscript contributes executable infrastructure: a shared scenario contract, a traceable formulation, and reproducible baseline evidence. Relative to companion modelling studies, it provides the stable software and data baseline used to test methodological questions (e.g., rolling-horizon configuration effects), so platform and modelling claims can be evaluated separately.

\subsection{Maturity and current scope}
The current release is an early public platform baseline. It demonstrates reproducible scheduling experiments on public BC reference scenarios and synthetic tiers (Section~\ref{sec:illustrative-example}), but does not claim broad external deployment evidence.

\subsection{Transferability and limits}
FHOPS is structured to support transfer to new contexts through scenario and parameter updates rather than solver rewrites. Three current limits remain explicit:
\begin{itemize}
  \item Validation breadth is still centered on the current public reference ladder.
  \item Runtime comparisons are environment-dependent and should be interpreted as within-study baselines.
  \item Warm-start MIP integration is implemented, but decisive performance gains at larger stress cases are not yet claimed.
\end{itemize}

These limits keep the manuscript's claims aligned with demonstrated evidence.
